\documentclass{Phyport}

\usepackage{hyperref}
\hypersetup{
    colorlinks=true,
}
% 选用宏（如有额外需要的宏包请在此添加）
\usepackage{float}       % 提供 [H] 选项
\usepackage{booktabs}    % 三线表
\usepackage{amsmath}     % 数学公式
\usepackage{caption}     % 标题设置
\usepackage{siunitx}     % SI 单位排版
 %选用宏

\exname{铁磁材料居里温度测量} %实验名称
\extable{4} %实验桌号
\instructor{陈飞燕} %指导教师
\class{电科2502} %班级
\name{管云哲} %姓名
\stuid{3250104519} %学号

\nyear{2026} %年
\nmonth{9} %月
\nday{29} %日
\nweekday{二} %星期几

\redate{} %如有实验补做，补做日期
\resitu{} %情况说明：

\begin{document}
\setcounter{page}{1}
\makecover

\section{预习报告（10分）}
（注：将已经写好的“物理实验预习报告”内容拷贝过来）

\subsection{实验综述（5分）}
（自述实验现象、实验原理和实验方法，包括必要的光路图、电路图、公式等。不超过500字。）

铁磁质内部相邻原子间存在强烈的交换耦合作用，使电子自旋磁矩在小区域内自发整齐排列，形成磁畴，因而在无外磁场时即具有自发磁化。温度升高时，热运动逐渐破坏磁畴的有序排列；当温度达到居里温度 $T_C$ 时，磁畴瓦解，材料由铁磁态转变为顺磁态，磁导率 $\mu_r$ 由远大于 1 骤降至接近 1。$T_C$ 仅由化学成分与晶体结构决定，与晶粒大小、取向、应力等结构因素几乎无关，故称结构不灵敏参数。

磁介质的磁化规律为
\begin{equation}
B=\mu_0(H+M)=(\chi_m+1)\mu_0H=\mu_r\mu_0H=\mu H
\end{equation}

本实验采用交流电桥法测量居里温度。将铁氧体样品插入 $L_2$ 线圈后，线圈电感 $L=\mu N^2A/l\propto\mu(T)$ 随温度变化，破坏电桥平衡，桥臂间电位差 $U_{CD}$ 随之改变。取 $R_a=R_b$，电桥平衡条件为
\begin{equation}
R_w=\frac{R_b}{R_a}R_1,\qquad L_2=\frac{R_b}{R_a}L_1
\end{equation}

当电感变化 $\Delta L_2$ 时，桥路输出
\begin{equation}
\Delta U_{CD}\propto\frac{\partial U_{CD}}{\partial L_2}\Delta L_2\propto\frac{\omega\Delta L_2}{(R_b+R_w+j\omega L_2)^2}
\end{equation}

可见 $U_{CD}$ 非线性但单调地反映 $\mu$ 随温度的变化。缓慢升温（或降温），逐点记录温度 $T$ 与电压 $U$，作 $U\sim T$ 曲线，其斜率最大处（拐点）对应的温度即为居里温度 $T_C$。

\begin{figure}[H]
    \centering
    \includegraphics[width=0.62\linewidth]{mu_T_curve.png}
    \caption{磁导率 $\mu$ 与磁场 $H$ 及温度 $T$ 的关系（右图突变点对应居里温度 $T_C$）}
    \label{fig:muT}
\end{figure}

\subsection{实验重点（3分）}
（简述本实验的学习重点，不超过100字。）

理解磁畴与交换耦合作用下铁磁—顺磁转变的微观机理；掌握 RL 交流电桥的平衡条件与输出电压随电感变化的规律；学会由 $U\sim T$ 曲线确定居里温度并分析升温速率、信号频率对测量结果的影响。

\subsection{实验难点（2分）}
（简述本实验的实现难点，不超过100字。）

调平电桥需将电压调至 $1.0$\,mV 以下，需耐心微调 $R_w$；样品与温度传感器之间存在热滞后，快速加热会使测得 $T_C$ 偏高，须保持动态平衡；居里点附近 $U$ 变化剧烈，应加密取点（0.5\,$^\circ$C）以免漏掉拐点。

\section{实验操作与原始数据（50分）}
（将有老师签名的“自备数据记录草稿纸”的扫描或手机拍摄图粘贴在下方，完整保留姓名，学号，教师签字和日期。按老师的要求拍摄实验过程图，附在此处）

\begin{figure}[H]
    \centering
    \includegraphics[width=0.52\linewidth]{实验数据.jpg}
    \caption{原始数据}
    \label{fig:rawdata}
\end{figure}

\begin{figure}[H]
    \centering
    \includegraphics[width=0.92\linewidth]{实验过程.jpg}
    \caption{实验过程：FB2020A 型居里温度测量仪面板}
    \label{fig:process}
\end{figure}

\subsection{实验装置与操作步骤}

实验装置为 FB2020A 型居里温度测量仪，由数控信号源、交流电压表、PID 温度控制仪、交流电桥与加热样品架组成。测量步骤如下。

\begin{enumerate}
    \item 按图接线，样品先不插入，将“加热选择”置于“断”位；频率调至 $1000$\,Hz，波形选正弦波，幅度调至约中间位置，电压表量程先置 $2000$\,mV 档。
    \item 调节 $R_w$ 使交流电压表读数最小，本次调至 $1.0$\,mV 以下，电桥平衡。
    \item 在铁氧体样品棒上均匀涂导热硅脂后插入加热架右侧空心线圈的铜管中，电桥失平，电压表读数升至 $100$\,mV 以上，据示值及时切换至 $200$\,mV 档以提高分辨率。
    \item “加热选择”置于“低”档，由 $40\,^\circ$C 起记录 $U$；在 $50$\,$^\circ$C 之后每 $5\,^\circ$C 记一次，至 $75\,^\circ$C 附近电压开始明显下降后改为每 $0.5\,^\circ$C 记录一次，直至 $120\,^\circ$C 附近关闭加热。
    \item 全程保持信号源幅度与样品在线圈中的位置不变。
\end{enumerate}

\subsection{原始数据记录}

室温 $23.5\,^\circ$C，信号频率 $1000$\,Hz，波形为正弦波，样品为铁氧体，居里温度参考值 $78\,^\circ$C。原始电压以 mV 直读，交流电压表 $200$\,mV 档最小分辨率 $0.1$\,mV，温度控制精度 $0.5\,^\circ$C。按记录纸原样整理于表1。

\begin{table}[H]
    \centering
    \caption{铁氧体样品交流电桥输出电压与加热温度关系（升温）}
    \label{tab:data}
    \small
    \begin{tabular}{cccccc}
        \toprule
        $T$ / $^\circ$C & $U$ / mV & $T$ / $^\circ$C & $U$ / mV & $T$ / $^\circ$C & $U$ / mV \\
        \midrule
        40.0 & 36.0 &  80.5 & 29.9 &  89.5 & 13.8 \\
        45.0 & 36.2 &  81.0 & 28.6 &  90.0 & 13.3 \\
        50.0 & 36.4 &  81.5 & 27.2 &  90.5 & 13.0 \\
        55.0 & 36.5 &  82.0 & 25.9 &  91.0 & 12.6 \\
        60.0 & 36.7 &  82.5 & 24.2 &  91.5 & 12.2 \\
        65.0 & 36.8 &  83.0 & 22.4 &  92.0 & 11.9 \\
        70.0 & 36.9 &  83.5 & 21.5 &  92.5 & 11.6 \\
        75.0 & 36.6 &  84.0 & 20.5 &  93.0 & 11.3 \\
        75.5 & 36.4 &  84.5 & 19.6 &  93.5 & 11.0 \\
        76.0 & 36.2 &  85.0 & 18.9 &  94.0 & 10.7 \\
        76.5 & 36.0 &  85.5 & 18.2 &  94.5 & 10.4 \\
        77.0 & 35.6 &  86.0 & 17.5 &  95.0 & 10.2 \\
        77.5 & 35.3 &  86.5 & 16.9 & 100.0 &  8.0 \\
        78.0 & 34.3 &  87.0 & 16.4 & 105.0 &  6.5 \\
        78.5 & 33.9 &  87.5 & 15.7 & 110.0 &  4.7 \\
        79.0 & 33.1 &  88.0 & 15.2 & 115.0 &  3.2 \\
        79.5 & 32.4 &  88.5 & 14.5 & 120.0 &  2.4 \\
        80.0 & 31.2 &  89.0 & 14.1 &       &      \\
        \bottomrule
    \end{tabular}
\end{table}

\section{结果与分析（30分）}
\subsection{数据处理与结果（20分）}
（列出数据表格、选择适合的数据处理方法、写出测量或计算结果。）

\subsubsection{升温 $U\sim T$ 曲线}

将表1数据以温度 $T$ 为横坐标、桥路输出电压 $U$ 为纵坐标作图，得图~\ref{fig:UT}。曲线可分为三段：$40$–$75\,^\circ$C 段 $U$ 基本平坦，由 $36.0$\,mV 缓慢升至极大值 $36.9$\,mV，此即铁磁态下磁导率随温度缓慢变化的反映；$75\,^\circ$C 之后 $U$ 迅速单调下降，$78$–$82\,^\circ$C 区间下降最陡，斜率绝对值达到最大；$95\,^\circ$C 以后曲线趋于平缓，$120\,^\circ$C 时 $U=2.4$\,mV，仍有残余电压而未回到零。

\begin{figure}[H]
    \centering
    \includegraphics[width=0.86\linewidth]{UT_heating.pdf}
    \caption{铁氧体样品升温 $U\sim T$ 曲线（散点为实测数据，实线为四参数光滑拟合，虚线为拐点法确定的居里温度）}
    \label{fig:UT}
\end{figure}

\subsubsection{拐点法确定居里温度}

对实测数据作三次样条插值得到光滑曲线 $U(T)$，再求一阶导数 $\mathrm{d}U/\mathrm{d}T$，其绝对值最大处即为居里温度。在 $70$–$95\,^\circ$C 区间计算得
\begin{equation}
\left.\frac{\mathrm{d}U}{\mathrm{d}T}\right|_{\min}=-2.657\ \mathrm{mV/^\circ C},
\qquad T_C=81.72\ ^\circ\mathrm{C}
\end{equation}

斜率曲线见图~\ref{fig:dudt}，其谷底即拐点位置。为进一步验证，另用最小二乘法将数据拟合为“顺磁基线 $+$ 相变台阶”的四参数模型
\begin{equation}
U(T)=U_0+A\left[1+k\frac{T-100}{20}\right]\frac{1}{1+\exp\left(\dfrac{T-T_C}{w}\right)}
\end{equation}
其中 $U_0$ 为顺磁基线、$A$ 为铁磁态富余电压、$w$ 为相变温度宽度、$k$ 为基线斜率修正。拟合得 $U_0=6.25$\,mV、$A=40.28$\,mV、$w=4.42\,^\circ$C、$T_C=82.87\,^\circ$C，残差均方根 $1.370$\,mV。两法所得 $T_C$ 相差 $1.15\,^\circ$C，取推荐度更高的样条拐点法结果为测量值。

\begin{figure}[H]
    \centering
    \includegraphics[width=0.86\linewidth]{dU_dT.pdf}
    \caption{$U\sim T$ 曲线的斜率曲线，其绝对值最大处对应居里温度}
    \label{fig:dudt}
\end{figure}

\subsubsection{不确定度评定}

居里温度的不确定度由 A、B 两类分量合成。

A 类分量来自拐点定位的稳定性。将参与拟合的温度区间分别取 $(70,92)$、$(70,95)$、$(70,98)$、$(68,95)$、$(72,95)$、$(72,98)$\,$^\circ$C 六种情形重复计算，得 $T_C=81.60,81.60,81.58,81.60,81.58,81.60\,^\circ$C，样本标准差
\begin{equation}
u_A=\sqrt{\frac{1}{n-1}\sum_{i=1}^{n}\left(T_{C,i}-\overline{T_C}\right)^2}=0.01\ ^\circ\mathrm{C}
\end{equation}

B 类分量来自仪器误差。PID 温控仪控温精度 $0.5\,^\circ$C，按均匀分布取 $u_{B,T}=0.5/\sqrt{3}=0.29\,^\circ$C；交流电压表在该量程下的分辨极限 $0.1$\,mV，经斜率折算为温度 $u_{B,U}=0.1/(\sqrt{3}\times 2.657)=0.022\,^\circ$C。二者合成为
\begin{equation}
u_B=\sqrt{u_{B,T}^2+u_{B,U}^2}=\sqrt{0.29^2+0.022^2}=0.29\ ^\circ\mathrm{C}
\end{equation}

合成标准不确定度为
\begin{equation}
u=\sqrt{u_A^2+u_B^2}=\sqrt{0.01^2+0.29^2}=0.29\ ^\circ\mathrm{C}
\end{equation}

因此测量结果为
\begin{equation}
T_C=(81.7\pm 0.3)\ ^\circ\mathrm{C}\qquad (k=1)
\end{equation}

与样品给出的参考值 $78\,^\circ$C 相比，相对偏差
\begin{equation}
E_r=\frac{\left|81.7-78\right|}{78}\times 100\%=-\;4.8\%
\end{equation}
即测得值比参考值偏高约 $3.7\,^\circ$C。

\subsection{误差分析（10分）}
（运用测量误差、相对误差或不确定度等分析实验结果，写出完整的结果表达式，并分析误差原因。）

最终结果表达式为 $T_C=(81.7\pm 0.3)\,^\circ$C $(k=1)$，相对不确定度 $u_r=u/T_C=0.37\%$。测得值偏高 $3.7\,^\circ$C，相对偏差 $4.7\%$，误差来源分析如下。

\textbf{（1）样品与传感器间的热滞后，这是主要的系统误差来源。}
温度传感器贴于加热铜管，铁氧体样品插在管中，二者之间隔着铜管壁与导热硅脂。升温过程中样品温度始终低于传感器读数，且该温差随升温速率增大而增大。由讲义与课件数据，$78\,^\circ$C 样品在 $1000$\,Hz 下低档加热与高档加热测得 $74.5\,^\circ$C 与 $79.0\,^\circ$C，差值达 $4.5\,^\circ$C，与本次偏高量级一致。本次采用“低”档加热并在居里区间加密取点，已将该效应压至较小，但加热速率终究不为零，滞后不能完全消除，测得 $T_C$ 必然偏高。这也是课件指出“对本实验装置，利用降温曲线测居里温度更加准确”的原因：降温时样品温度高于传感器读数，偏差方向相反，可用升、降温两条曲线取平均以抵消热滞后。本次未记录降温数据，无法作此修正，是偏高的首要原因。

\textbf{（2）电桥未完全调平引入的偏置。}
预备实验中电桥调平时电压需降至 $1.0$\,mV 以下，实际难以精确到零。电桥残余不平衡电压 $U_0$ 叠加在信号上，相当于把整条 $U\sim T$ 曲线整体平移。课件数据显示，调平电压由 $1.3$\,mV 变为 $10.0$\,mV 时，升温曲线测得 $T_C$ 由 $74.5\,^\circ$C 变为 $77.5\,^\circ$C，可见影响显著。本次已将调平电压控制在 $1.0$\,mV 以下，该项贡献估计不大于 $0.5\,^\circ$C，属次要因素。

\textbf{（3）相变温度本身的分布与采样间隔。}
实际铁氧体存在成分起伏与晶界、位错等缺陷，各微区居里点并不相同，而是在 $w=4.42\,^\circ$C 的温区内分布，居里温度本身并非一个锐利单值。不同的 $T_C$ 定义方式（拐点法、切线外推法、区间上下限法）将给出不同的数值，本次两法相差 $1.15\,^\circ$C 即源于此，这是方法本身固有的不确定度而非操作失误。

\textbf{（4）其他次要误差。}
交流电压表读数引入 $\pm 0.1$\,mV 的随机误差，经斜率折算为 $0.02\,^\circ$C，可忽略；温度记录以 $0.5\,^\circ$C 为间隔，拐点定位被量化误差限制在 $\pm 0.25\,^\circ$C 以内，已计入 A 类分量；此外样品棒插入深度、与线圈的相对位置在每次装夹时的微小差异会改变耦合系数，进而影响 $U$ 的绝对大小。实验中保持样品位置不变是必要的，本次已严格遵守。

综上，本实验结果的偏差属以热滞后为主导的系统误差，而非随机误差。减小误差的根本途径是改用降温曲线测量并与升温曲线取平均，或将升温速率进一步降低至准静态。

\section{实验探讨与思考题（10分）}
（对实验内容、现象和过程的小结，不超过100字。并解答教材或讲义或老师布置的思考题，请先写题干，再作答。）

本实验以 RL 交流电桥将铁氧体磁导率的温度变化转换为桥路电压，由 $U\sim T$ 曲线的拐点定出 $T_C=(81.7\pm 0.3)\,^\circ$C，偏高参考值 $4.7\%$。体会最深的是：热滞后与电桥残余偏置都是方向确定的系统误差，只有改变测量条件（降温、减小加热速率）才能消除，靠增加测点数量无济于事。

\textbf{思考题 1：铁磁物质的三个特性是什么？}

铁磁性物质具有磁畴结构，与之相应表现出三个基本特性：一是\textbf{高磁导率}，其相对磁导率可达 $10^2$–$10^5$ 量级，远大于顺磁质的 $1+\chi_m\sim 1+10^{-5}$；二是\textbf{磁滞现象}，磁化状态与磁化历史有关，外场去掉后存在剩磁，$B\sim H$ 曲线形成回线，源于掺杂与内应力阻碍磁畴复原；三是存在\textbf{居里温度}，超过该温度磁畴瓦解，高磁导率与磁滞等铁磁特性全部消失，材料转为顺磁质。三者都与磁畴结构的存在直接相关。

\textbf{思考题 2：用磁畴理论解释样品的磁化强度在温度达到居里点时发生突变的微观机理是什么？}

铁磁质内部存在两种相互竞争的效应：相邻原子间的交换耦合作用使电子自旋磁矩倾向于平行排列，是磁畴自发磁化的根源；热运动则使自旋取向随机化。低温下交换耦合能量远大于 $k_BT$，自旋整齐排列，磁畴稳定存在，饱和磁化强度接近其零温值。温度升高时热扰动增强，自旋取向逐渐偏离有序状态，饱和自发磁化强度按布洛赫 $T^{3/2}$ 律缓慢下降。当温度升至 $T_C$ 时，热运动能量与交换作用能量相当，交换作用再也无法维持自旋的平行排列，磁畴随之瓦解，自发磁化强度在很窄的温区内急剧降至零，材料由铁磁态转变为顺磁态。$T_C$ 由交换积分 $J$ 的强度决定，而 $J$ 取决于化学成分与晶体结构，故 $T_C$ 是结构不灵敏参数。

\textbf{思考题 3：测出的 $U\sim T$ 曲线，为什么与横坐标没有交点？}

理论上温度远高于 $T_C$ 后材料完全转为顺磁态，$\mu_r\to 1$，电感 $L_2$ 与插入样品前相差无几，电桥本应重新平衡、$U\to 0$，曲线似应与横坐标相交。但实测 $120\,^\circ$C 时仍有 $U=2.4$\,mV，且拟合给出的顺磁基线为 $U_0=6.25$\,mV，曲线整体悬在横坐标之上。原因有三：其一，预备实验时电桥的调平是“无样品”状态下的平衡，而测量时样品插入线圈后其自身的抗磁性背景与几何占位仍会改变等效电感，使顺磁态的平衡点与预备平衡点并不重合，存在恒定残余电压；其二，电桥调平本身只能到 $1.0$\,mV 量级，这个残余不平衡电压直接平移到全曲线上；其三，样品棒并非理想细长，其退磁因子随温度变化引起的微小电感变化在高阻温区仍有贡献。此外，$120\,^\circ$C 时材料虽已过 $T_C$，但自发磁化并非在 $T_C$ 处阶跃归零，而是沿 $\left(1-T/T_C\right)^{1/2}$ 型的临界规律递减，在 $120\,^\circ$C 仍残留极弱磁化。故曲线是逐渐逼近横坐标而非与它相交。

{\fangsong 注意事项：
\begin{enumerate}[label=\arabic*.]
    \item 用PDF格式上传“实验报告”，文件名：学生姓名+学号+实验名称+周次。
    \item “实验报告”必须递交在“学在浙大”本课程内对应实验项目的“作业”模块内。
    \item “实验报告”成绩必须在“浙江大学物理实验教学中心网站”-“选课系统”内查询。
    \item 教学评价必须在“浙江大学物理实验教学中心网站”-“选课系统”内进行，学生必须进行教学评价，才能看到实验报告成绩，教学评价须在本次实验结束后3天内进行。
\end{enumerate}}
\begin{center}
    \fbfangsong{浙江大学物理实验教学中心制}
\end{center}
\end{document}
